% Source: https://artofproblemsolving.com/community/c3h2024484_geometeryishcounting_question
% Stable identity: aops-topic-2024484-q01
\textbf{Upper bound.} Take two of the polygons, with $m<n$ sides. Each side of the $m$-gon is a segment not lying on any side line of the convex $n$-gon, so it meets the boundary of the $n$-gon in at most $2$ points (a line meets the boundary of a convex polygon in at most two points unless it contains a side). Hence these two polygons meet in at most $2m=2\min(m,n)$ points. Since no point lies on three polygons, each intersection point belongs to exactly one pair of polygons, and the total is at most
\[
\sum_{\text{pairs } m<n}2m=\sum_{m=3}^{49}2m(50-m),
\]
because for each $m$ there are $50-m$ polygons with more sides.

\textbf{Construction.} Inscribe regular polygons with $3,4,\dots,50$ sides in one circle, rotated so that no two share a vertex and no three sides pass through a common point (all but finitely many choices of rotation angles work). Consider an $m$-gon and an $n$-gon with $m<n$. Each side of the $m$-gon cuts off an arc of $\dfrac{360^\circ}{m}>\dfrac{360^\circ}{n}$, which contains at least one vertex of the $n$-gon but not all of them; these vertices are consecutive around the circle. Two chords of a circle cross exactly when their endpoints alternate, so exactly two sides of the $n$-gon cross this side of the $m$-gon (the ones entering and leaving that arc). Thus the two polygons meet in exactly $2m$ points, and the bound is attained. The figure shows a triangle and a square meeting in $6$ points.

\begin{asy}
usepackage("amsmath");
size(4.2cm);
defaultpen(fontsize(10pt));
draw(unitcircle, gray+dashed);
pair[] T, Q;
for(int i=0;i<3;++i) T.push(dir(90+120*i));
for(int i=0;i<4;++i) Q.push(dir(20+90*i));
draw(T[0]--T[1]--T[2]--cycle, blue+linewidth(0.9));
draw(Q[0]--Q[1]--Q[2]--Q[3]--cycle, red+linewidth(0.9));
for(int i=0;i<3;++i) for(int j=0;j<4;++j){
  pair a=T[i], bb=T[(i+1)%3], c=Q[j], d=Q[(j+1)%4];
  real den=cross(bb-a,d-c);
  if(abs(den)>1e-9){
    real t=cross(c-a,d-c)/den, u=cross(c-a,bb-a)/den;
    if(t>0 && t<1 && u>0 && u<1) dot(a+t*(bb-a));
  }
}
label("\mbox{}",(min(currentpicture,user=true)+max(currentpicture,user=true))/2,black); // keeps the text after the figure black when LaTeX typesets the labels inline
\end{asy}

\textbf{Computation.} Using $\sum_{m=3}^{49}m=\dfrac{49\cdot50}{2}-3=1222$ and $\sum_{m=3}^{49}m^2=\dfrac{49\cdot50\cdot99}{6}-5=40420$,
\[
\sum_{m=3}^{49}2m(50-m)=100\cdot1222-2\cdot40420=122200-80840=\boxed{41360}.
\]
